\subsection{Logique de décision et interprétation des commandes}

Après le traitement du langage naturel par l'API Gemini, la réponse est analysée pour déterminer le type d'action à exécuter.

\begin{table}[H]
\centering
\caption{Correspondance entre commandes vocales et actions}
\begin{tabular}{|c|c|}
\hline
\textbf{Commande vocale} & \textbf{Action exécutée} \\
\hline
``ouvre la bouche'' & Ouverture mâchoire \\
``ferme la bouche'' & Fermeture mâchoire \\
``regarde à gauche'' & Yeux à gauche \\
``regarde à droite'' & Yeux à droite \\
``fais une expression joyeuse'' & Expression faciale \\
Question & Réponse vocale \\
\hline
\end{tabular}
\end{table}

\begin{verbatim}
if intent == "ouvrir_bouche":
    publier("/axel/intent_action", IntentAction(
        action_type="jaw", function_name="ouvrir_bouche"))
elif intent == "fermer_bouche":
    publier("/axel/intent_action", IntentAction(
        action_type="jaw", function_name="fermer_bouche"))
elif intent == "regarder_gauche":
    publier("/axel/intent_action", IntentAction(
        action_type="eyes", function_name="regarder_gauche"))
elif intent == "regarder_droite":
    publier("/axel/intent_action", IntentAction(
        action_type="eyes", function_name="regarder_droite"))
elif intent == "expression_joie":
    publier("/axel/intent_action", IntentAction(
        action_type="expression", function_name="expression_joie"))
else:
    publier("/audio/output", "réponse vocale")
\end{verbatim}
